% Generated draft: requires mathematical and source review.
\documentclass[11pt]{article}
\usepackage[margin=25mm]{geometry}
\usepackage{fontspec}
\setmainfont{Latin Modern Roman}
\usepackage{amsmath,amssymb,mathtools,unicode-math}
\setmathfont{Latin Modern Math}
\usepackage{enumitem,needspace,xurl,hyperref}
\setlength{\parindent}{0pt}
\setlength{\parskip}{5pt}
\setlength{\emergencystretch}{3em}
\newcommand{\N}{\mathbb{N}}
\newcommand{\Z}{\mathbb{Z}}
\newcommand{\Q}{\mathbb{Q}}
\newcommand{\R}{\mathbb{R}}
\newcommand{\C}{\mathbb{C}}
\newcommand{\F}{\mathbb{F}}
\newcommand{\A}{\mathbb{A}}
\newcommand{\PP}{\mathbb{P}}
\newcommand{\E}{\mathbb{E}}
\newcommand{\eps}{\varepsilon}
\newcommand{\dd}{\,\mathrm{d}}
\begin{document}
\section*{UnsolvedMath ID: 701239}
\textbf{Source problem code:} catalogue entry
\subsection{Definitions and mathematical statement}
Write an abelian group \(G\) additively, with identity \(0\). It is torsion-free if, for every \(g\in G\) and every integer \(n\geq 1\), the equality \(ng=0\) implies \(g=0\). A lattice ordering on \(G\) is a partial order \(\leq\) such that every pair \(x,y\in G\) has a greatest lower bound \(x\wedge y\) and a least upper bound \(x\vee y\), and such that
\[
 x\leq y\ \Longrightarrow\ x+z\leq y+z
 \qquad(x,y,z\in G).
\]
Such an ordering is Archimedean if, for all \(x,y\in G\) with \(0\leq x\) and \(0\leq y\),
\[
 \bigl(nx\leq y\text{ for every integer }n\geq 1\bigr)
 \ \Longrightarrow\ x=0.
\]
A lattice ordering need not be a total order.

\textbf{Question (Kourovka Notebook 16.47).} Is it true that for every torsion-free abelian group \(G\), there exists an Archimedean lattice ordering on \(G\)? No restriction on the cardinality of \(G\) is stated.

The supplied statement is prefaced by ``(P. Conrad)'', while the record lists V. M. Kopytov and N. Ya. Medvedev as proposers; the respective attribution roles are not further explained. The supplied registry labels the problem open; this status is not independently verified here.
\subsection{Short English statement}
Does every torsion-free abelian group admit a translation-invariant partial order in which every pair has a meet and a join, and no nonzero positive element has all its positive integer multiples bounded above?
\subsection{Sources}
\begin{itemize}
\item Local UnsolvedMath registry, supplied complete source record: permanent ID \(701239\), exact source problem code \texttt{catalogue entry}, title \emph{Kourovka Notebook Problem 16.47}. The supplied status is \texttt{open}.
\item \emph{The Kourovka Notebook (Unsolved Problems in Group Theory)}, 16th issue (2006), Problem 16.47; also identified in the supplied background with the 21st issue edition (2026), \texttt{arXiv:1401.0300v48}. These bibliographic details are supplied source data and have not been independently checked.
\item Attribution supplied by the record: V. M. Kopytov and N. Ya. Medvedev, proposed year 2006. The statement additionally carries the parenthetical attribution ``(P. Conrad)''.
\item The supplied collection description names E. I. Khukhro and V. D. Mazurov as editors of the Kourovka Notebook collection covering Issues 1--21 (1965--2026).
\end{itemize}
\end{document}
